\chapter{Concord}

% F009, F029
A grape is a berry. This is not a figure of speech. It is a botanical
classification, and, like many botanical classifications, it has the effect
of making the ordinary seem improbable and the improbable seem
ordinary.\note{A strawberry is not a berry. It is an accessory fruit: most
of its flesh grows from the receptacle, the swollen tip of the stem, and
the tiny specks on its surface, each one an achene, are the actual fruits.
A raspberry is not a berry either. It is an aggregate fruit, a cluster of
drupelets, each with its own seed, its own thin skin. The grape, by
contrast, is the real thing: skin, flesh, and seeds, all from the ovary
wall, no stone.} Few things are entirely what they seem. The grape is a
berry, and a strawberry is not, and anyone who has trouble with this
distinction has not spent enough time around plants.

% F009
The skin of a grape is the exocarp. The flesh is the mesocarp. The thin
layers around the seeds are the endocarp. The ovary has two compartments,
each with two ovules, so a grape can have at most four seeds, and seeded
cultivars have one to four. The waxy film on the outside of the skin is the
bloom. It is the thing you see on a Concord grape that makes it look as
though it has just come in from the cold. Wipe it off and the grape turns
darker, shinier, more obvious, and less itself.

% F003, F004, F005
The grape belongs to the genus \textit{Vitis}, in the family
Vitaceae.\note[vitaceae]{The family has sixteen genera and about nine
hundred and fifty species, if you follow Wen, and twelve genera if you
follow Britannica, and neither is wrong, exactly; they are counting
different things in different ways, which is what taxonomists
do.\subnote{Virginia creeper and Boston ivy belong to the same family.
% F028
Most of its members climb.}} \textit{Vitis} has sixty to eighty species,
depending on whom you ask, and the reason no one can agree is that the
Asian species are poorly defined, which is a polite way of saying that not
enough botanists have gone to look at them. Most species are native to two
places: eastern Asia and North America. A single species,
\textit{vinifera}, is European and West Asian, and it is the one that makes
nearly all the wine in the world.

% F006
The wild ancestor of the cultivated grape is \textit{Vitis vinifera}
subspecies \textit{sylvestris}, which has male and female plants. The
cultivated vine has hermaphrodite flowers, which simplifies things
considerably if you are a vine. Domestication happened in western Asia,
about eight thousand years ago. A long time to have been at it, and still
not to agree on how many varieties there are.

% F007
\anchor{varieties}The OIV---the International Organisation of Vine and
Wine, which is based in Paris, where organisations of vine and wine tend to
be based---speaks of ten thousand known grapevine varieties. The Vitis
International Variety Catalogue lists twenty-one thousand and forty-five
variety names, of which twelve thousand two hundred and fifty are
\textit{vinifera}, but many of these are synonyms, the same grape wearing
different names in different countries.\note[names]{A grape that travels
collects names, as any traveller does.} The true number of distinct
\textit{vinifera} varieties is estimated at about six thousand. Of these,
thirteen cover more than a third of the world's vineyard area, and
thirty-three cover half. The grape has been generous with its variations
but stingy with its success.

% F008, F024
In 2018 the world produced seventy-seven point eight million tonnes of
grapes.\note{Fifty-seven per cent became wine, must, or juice. Thirty-six
per cent were table grapes. Seven per cent were dried---and it takes
roughly four kilograms of fresh grapes to make one kilogram of raisins, so
the raisin is a kind of concentrate.\subnote{The largest producers, in
millions of tonnes: China, 11.7; Italy, 8.6; the United States, 6.9;
Spain, 6.9; France, 6.2. The world's population that year was 7.7 billion.
Divide it out and you get about ten kilograms per person, or twenty-two
pounds, roughly the weight of a small child's bicycle.}} Whether any given
person actually received twenty-two pounds of grapes is another question.
Distribution is not the same as division.

% F010, F011
A bunch of grapes is botanically a panicle.\note{Its stalk before the
first branch is the peduncle, its main axis the rachis; each berry hangs
on a pedicel. A cluster can carry as many as a thousand flowers; fruit set
is considered good if about one flower in three becomes a berry.} I bought
a bunch in January. They were red, seedless, from Chile, which is where
% F023
American table grapes come from in the winter, because in the southern
hemisphere it is summer, and a grape does not know what country it is in.
The label said ``Product of Chile.'' The bunch sat on my desk for three
days before I counted it. I do not know why I counted it. I have been
paid, for most of my life, to be careful, and counting is a form of care,
or a form of habit, which may be the same thing.

The first count was a hundred and seventeen. I counted again and got a
hundred and eighteen. I counted a third time, slowly, moving each berry
aside with the tip of a pencil, and got a hundred and seventeen. The third
count agreed with the first, which ought to have settled the matter, but it
did not settle it entirely, because I do not know what I miscounted in the
second pass, and a discrepancy, even a discrepancy of one, will stay in
the mind the way a wrong note stays in the ear long after the orchestra
has moved on.

% F001
The first grape I knew was not, strictly speaking, a grape. It was Welch's
grape jelly, in a jar that became a glass. In the early
nineteen-sixties, Welch's sold its jelly in tumblers printed with
Flintstones characters, about four and a quarter inches tall, and when the
jelly was gone the glass went into the cupboard and stayed there, because
nobody in a small town in Iowa threw away a glass.\note{I do not remember
which character was on mine. I remember the jelly.} I remember it on white
bread with peanut butter, which is the combination that Americans call a
peanut-butter-and-jelly sandwich, as though there were no other jelly and
no other sandwich, which, in the rural Midwest in 1960, there more or less
was not.

% F020
\note[pbj]{The first published recipe for the sandwich appeared in 1901,
in the \textit{Boston Cooking-School Magazine}, by Julia Davis Chandler,
and the jelly it called for was currant or crab-apple, not grape. Grape
came later.\subnote{The sandwich became a childhood staple during the
Depression, helped by sliced bread, which arrived in 1928, and cheaper
peanut butter. The claim that it began with soldiers in the Second World
War is folklore, though soldiers did eat it.}} By the time I was eating it
the substitution was so complete that the word ``jelly,'' unmodified, meant
grape, the way the word ``Coke,'' in certain parts of the South, means any
soft drink at all. We are loyal to the first few things we are given, and
we call the loyalty taste.

% F013
The taste that Americans call ``grape'' is the taste of one
variety.\note[methyl]{The compound responsible is methyl anthranilate,
first identified in grape juice in 1921, by Power and Chesnut. It is found
in \textit{labrusca}-type grapes but not in \textit{vinifera}. Other
compounds contribute---furaneol, 2-aminoacetophenone---but methyl
anthranilate is the signature.\subnote[foxy]{The French have a word for the
flavour: \textit{fox\'e}, from \textit{go\^ut fox\'e}, the foxy taste,
% F027
which comes from the English name for the species, the fox grape. Whether
foxes eat the grape, or the grape smells like a fox, is a question I have
not been able to resolve.}} It is the taste of the Concord. The distance
between a Concord and a Chardonnay is the distance between the New World
and the Old, which is a geographical fact before it is anything else.

% F014, F015
The Concord was made by a man named Ephraim Wales Bull, who was born in
Boston on March 4, 1806, and who moved to Concord, Massachusetts, in
August of 1836, to a house on Lexington Road that is now called Grapevine
Cottage. In the autumn of 1843, by his own account, he sowed seeds from a
wild grape, and by 1849 the selected vine bore good fruit. He exhibited it
at the Massachusetts Horticultural Society in 1853. It was first sold by
C.\,M.~Hovey, of Boston, in 1854.

% F017
\anchor{concord}The Concord is usually called a cultivar of \textit{Vitis
labrusca}, but it is a hybrid, with some \textit{vinifera} in its
ancestry.\note{Huber and others, in 2016, identified it as an offspring of
`Catawba' and an unknown wild parent.\subnote{Few things are entirely what
they seem. The most American of grapes is not entirely American.}}

% F025
Bull wanted, in his own words, ``a good table and wine grape, which shall
also be early, hardy and prolific.'' He got all three. The grapes then
available in New England died in the winter or could not ripen before the
first autumn frost. The European \textit{vinifera} vines are not
cold-hardy and are prone to disease in the eastern United States. The
Concord ripened early and survived hard winters. It was exactly what was
needed, and Bull gave it to the world, and the world took it.

% F016
Bull served in the Massachusetts legislature. He died in 1895 at the
Concord Home for the Aged. He is buried in Sleepy Hollow Cemetery, and his
epitaph reads, ``He sowed, others reaped.''\note[epitaph]{The words are
widely quoted. The punctuation varies between sources. I have not seen the
stone. I mean to.}

% F018
The man who reaped was a dentist. Thomas Bramwell Welch was born in
Glastonbury, England, in 1825. He came to the United States, became a
Methodist, became a prohibitionist, and settled in Vineland, New Jersey,
where he practised dentistry and disapproved of alcohol. In 1869 he took
forty pounds of Concord grapes---Welch's says forty pounds---from in front
of his house and pasteurized the juice, which kept it from fermenting,
which kept it from becoming wine, which was the point. He wanted an
unfermented grape juice for communion. The Lord's Supper without the
fermentation. He succeeded.

% F019, F026
His son, Charles, turned the product into a
business.\note{In 1897 the business built a plant in Westfield, New York,
on Lake Erie, in the heart of the Concord Grape Belt, a strip about sixty
miles long through Chautauqua County. Welch's sold Grapelade, a grape
spread, from 1918, and grape jelly from 1923, and the jelly was made from
Concord grapes, and it still is.} It is the jelly that I ate as a child in
Iowa, on white bread, with peanut butter, in a kitchen I can see but
cannot return to, in a town I will not name, because someone from there
will write to correct me, and I would rather be uncorrected and free.

We honour the man who profits and forget the one who planted. This is not a
thing I say with bitterness. It is a thing I say with interest, the way
you might note which direction the wind is blowing without wishing it to
change. Bull sowed. Welch reaped. The grape is indifferent to both of
them.

% F002
The bunch is still on my desk. It is late January. Outside, the vines are
dormant, and the pruning has not yet begun. In Napa the growers wait until
February or March, because January rain carries fungal spores into fresh
cuts, and a fresh cut on a vine is an invitation that should not be
extended lightly. January is for marking the wood---deciding what stays and
what goes, which is a kind of editing. You look at what grew last year. You
decide what will be allowed to grow next year. You cut everything else. If
you are too generous the vine makes too many clusters and none of them
ripen well. If you are too severe you lose a year. The vine does not tell
you which you have been. You find out in September.

I will count the grapes again in the morning. I expect a hundred and six.
